\documentclass[a4paper,10pt]{article}
\usepackage[utf8]{inputenc}

\title{Especificaciones soportadas en el QFTbx}
\author{Isaac Martínez Forte}

\begin{document}

\maketitle

Con $P$ una planta cualquiera de la familia, $C$ el controlador,
$L = PC$ el lazo de esa planta y $F$ el prefiltro, las seis restricciones
que QFTbx admite son las que recoge la literatura (tesis, 2022,
ecuaciones 1.6 a 1.11). Todas son cotas sobre la magnitud de una función de
transferencia en lazo cerrado, y todas han de cumplirse para toda la familia
de plantas.

Seguimiento:

$$\alpha(\omega) \leq \left| F\frac{L}{1+L} \right| \leq \beta(\omega)$$

Estabilidad robusta:

$$\left| \frac{L}{1+L} \right| \leq \lambda$$

Eliminación de ruido del sensor:

$$\left| \frac{L}{1+L} \right| \leq \delta_n(\omega)$$

Rechazo de perturbaciones a la salida de la planta (RPS):

$$\left| \frac{1}{1+L} \right| \leq \delta_{po}(\omega)$$

Rechazo de perturbaciones a la entrada de la planta (RPE):

$$\left| \frac{P}{1+L} \right| \leq \delta_{pi}(\omega)$$

Esfuerzo de control:

$$\left| \frac{C}{1+L} \right| \leq \delta_{ce}(\omega)$$

\bigskip

Dos observaciones sobre lo que el programa hace con ellas. El prefiltro $F$
multiplica el lazo cerrado entero en cada frecuencia, de modo que desplaza la
banda de $|FL/(1+L)|$ pero no la estrecha: lo que el ajuste del lazo tiene que
conseguir es que la DISPERSIÓN de $|L/(1+L)|$ sobre la familia quepa dentro de
$\beta - \alpha$, y el prefiltro se diseña después para deslizarla a su sitio.
QFTbx calcula la frontera de seguimiento contra esa diferencia y todavía no
diseña $F$. Y no hay sensor $H$: un sensor distinto de 1 forma parte de $P$.

\end{document}
